\chapter{物理架构}

\section{宿主与内核}

目标实现由 TypeScript 宿主和 Rust 内核组成。TypeScript 宿主负责模型服务、外部资源适配、调用方上下文、离线认知工作流、维护租约集成和面向产品的编排；Rust 内核负责认知领域语义、权威写入验证、持久化一致性、认知查询、检索计划、服务投影生命周期和管理操作。

宿主不得直接绕过内核修改认知权威。内核不直接拥有模型供应商账户、外部资源凭据或产品行为调度。

\begin{ReportWideFigure}
\centering
\begin{tikzpicture}[
  block/.style={draw=NWBlue,rounded corners,align=center,text width=5.0cm,minimum height=1.15cm,inner sep=6pt},
  data/.style={draw=NWRule,rounded corners,fill=NWLight,align=center,text width=3.4cm,minimum height=0.9cm,inner sep=5pt},
  arrow/.style={-{Stealth[length=2mm]},thick,color=NWBlue}
]
\node[block] (consumer) {调用方 / Heptalogos\\认知请求、上下文贡献与结果回执};
\node[block,below=0.65cm of consumer] (host) {Nous Wave 宿主（TypeScript）\\模型、资源、上下文与离线流程};
\node[block,below=0.65cm of host] (kernel) {Nous Wave 内核（Rust）\\领域所有者、权威数据、查询与服务生命周期};
\node[data,below left=0.85cm and 0.4cm of kernel] (pg) {PostgreSQL\\结构化权威数据};
\node[data,below=0.85cm of kernel] (lex) {词项检索\\Tantivy};
\node[data,below right=0.85cm and 0.4cm of kernel] (dense) {稠密检索\\可替换后端};
\node[data,right=0.9cm of host] (object) {不可变对象库\\OpenDAL 接入};
\draw[arrow] (consumer) -- (host);
\draw[arrow] (host) -- (kernel);
\draw[arrow] (kernel) -- (pg);
\draw[arrow] (kernel) -- (lex);
\draw[arrow] (kernel) -- (dense);
\draw[arrow] (host.east) -- (object.west);
\end{tikzpicture}
\captionof{figure}{Nous Wave 的目标物理架构。权威数据与可重建检索服务分别维护。}
\label{fig:physical-architecture}
\end{ReportWideFigure}

\section{持久化与检索组件}

PostgreSQL 保存结构化认知权威数据和事务状态。不可变对象库保存大型材料和二进制对象。Tantivy 提供词项检索；稠密检索后端和联想拓扑属于可替换的服务组件。所有服务组件都可以从权威数据和不可变材料重建。

稠密检索后端不在本设计中锁定具体产品。后端选择由本机工作负载、过滤能力、多视图支持、资源占用、备份恢复和运维成本的下游评测决定。

\section{部署原则}

参考部署以个人计算机本地运行优先，允许本地模型与远程模型服务并存。配置、数据和可执行实例分别维护，使实例升级或替换不要求删除主体数据。

管理接口为未来 CLI、管理 API 和图形界面提供稳定能力边界；目标设计不要求在 Nous Wave 仓库中实现特定 GUI。物理模块可以随着工程需要重组，只要不改变本文规定的语义所有权和系统边界。
